\documentclass[addpoints,12pt,A4]{exam}

\printanswers

\title{Continuous Optimization: Problem Set 3}

\input{ExSetup}


\begin{document}
\maketitle

\begin{questions}

\question Recall the definition of $\alpha$-self-concordance and barrier parameter $\theta$:
\[ \forall u,v,w: | D^{3} \phi(x) [u,v,w] | \leq \frac{2}{\sqrt{\alpha}} \|u \|_{x} \|v\|_{x} \|w\|_{x}   , \qquad \theta := \sup_{x \in D} \|(\nabla^{2}\phi(x) )^{-1}\nabla \phi(x)\|_{x} .  \]
\begin{enumerate}
    \item Assume $\phi_{k}$ are each $\alpha_{k}$-self-concordant with barrier parameter $\theta_{k}$ for domain $D$. Compute the self-concordance and barrier parameter for $\sum_{k} \phi_{k}$. 
    \item Show $- \log \det(X)$ is self-concordant for $D := \{ X \succ 0 \}$, and compute its barrier parameter. 
\end{enumerate}
\begin{solution}
    
\end{solution}



\question (Theorem 4.1.3 in Nesterov) We made the assumption that $\nabla^{2} \succ 0$ everywhere. In this question we remove the assumption. Show if there is some direction such that $\langle v, \nabla^{2} \phi(x) v \rangle = 0$ and $\phi$ is \emph{self-concordant}, then in fact 
\[ \forall y \in \text{dom}(\phi) : \langle v, \nabla^{2} \phi(y) v \rangle = 0 .   \]
What does this say about the domain in this direction?
\begin{solution}
    
\end{solution}



\question 
\begin{enumerate}
    \item Prove affine invariance of the Newton step: for convex $f$ let $g(y) := f(A y)$ so for $y = A^{-1} x$ $g(y) = f(x)$. Compute the gradient and Hessian and Newton step and show $y_{t+1} = A^{-1} x_{t+1}$ if $y_{t} = A^{-1} x_{t}$. 
    \item Prove affine invariance of the definition of self-concordance and barrier parameter. 
\end{enumerate}
\begin{solution}
    \begin{enumerate}
        \item 
    \end{enumerate}
\end{solution}



\question (Exercise 9.10 in BV): Newton's method with fixed step size $\lambda=1$ can diverge if the
initial point is not close to $x^{*}$. In this problem we consider two examples. 
\begin{itemize}
\item[a.] $f(x) = \log(e^{x} + e^{-x})$ has unique minimizer $x^{*} = 0$. Run Newton's method with
fixed step size $t=1$, starting at $x^{0} = 1$ and at $x^{0} = 1.1$. 


\item[b.] $f(x) = - \log(x) + x$ has unique minimizer $x^{*} = 1$. Run Newton's method with fixed
step size $\lambda=1$, starting at $x^{0} = 3$. 
\end{itemize}
\begin{solution}
    \begin{itemize}
        \item 
    \end{itemize}
    
\end{solution}



\question (Daniel Dadush Course, Ex 9.11)
Let $\MA \in \R^{m \times n}$, ${\rm rank}(\MA)=m$, $\vc \in \R^n$, $\vb \in \R^m$. 
For $\mu > 0$, examine the optimization problems
\begin{align}
\tag{BP} \label{barrier-primal}
\min_{\vx} &\{\vc \Tr \vx - \mu \sum_{i=1}^n \ln x_i: \MA \vx = \vb, \vx > 0 \} \\
\tag{BD} \label{barrier-dual}
\max_{\vs,\vy} &\{\vb \Tr \vy + \mu \sum_{i=1}^n \ln s_i: \MA \Tr \vy + \vs = \vc, \vs > 0\}. 
\end{align}
Assume that the feasible regions for both programs are non-empty (i.e., the
primal and dual LPs are both strictly feasible).
\begin{enumerate}
\item Show that the Lagrangian dual of~\ref{barrier-primal} is equivalent
to~\ref{barrier-dual} and prove that strong duality holds.
\item Conclude that $\vx_\mu$ is an optimal to solution to~\eqref{barrier-primal}
iff $\exists$ a solution $\vs_\mu, \vy_\mu$ to~\eqref{barrier-dual}
satisfying $x_{\mu,i} s_{\mu,i} = \mu$, $\forall i \in [n]$.
\end{enumerate}
\begin{solution}
    
\end{solution}


    
\end{questions}


\end{document}
